% a new subsection within the Background section)
%======================================================================

\subsection{GPS Ionospheric Delay Theory}
\label{sec:iono_theory}

The ionospheric height bias observed in the SD GPS positions (Section~\ref{sec:frame_problem}) is a well-understood consequence of operating single-frequency GPS receivers without an ionospheric correction model. This subsection provides the theoretical background.

\subsubsection{The Ionosphere as a Dispersive Medium}

The ionosphere is the region of Earth's upper atmosphere (approximately 60--1000\,km altitude) where solar ultraviolet radiation ionizes a sufficient fraction of the neutral gas to affect radio wave propagation. For a radio signal at frequency $f$ traversing the ionosphere, the refractive index for the group velocity (which governs code/pseudorange measurements) and the phase velocity (which governs carrier-phase measurements) differ from unity by opposite signs~\cite{misra2006,hofmann2008}:
\begin{align}
n_{\text{group}} &\approx 1 + \frac{40.3\,n_e}{f^2}, \label{eq:n_group}\\[4pt]
n_{\text{phase}} &\approx 1 - \frac{40.3\,n_e}{f^2}, \label{eq:n_phase}
\end{align}
where $n_e$ is the local electron density in electrons/m$^3$ and $f$ is in Hz. The group velocity is \emph{slowed} (the code pseudorange is too long), while the phase velocity is \emph{advanced} (the carrier phase appears too short). Both effects have the same magnitude.

Integrating the excess path length along the signal trajectory from satellite to receiver, the ionospheric group delay at frequency $f$ is:
\begin{equation}
\Delta I = \frac{40.3}{f^2}\int n_e\,ds = \frac{40.3}{f^2}\,\text{STEC},
\label{eq:iono_delay}
\end{equation}
where $\text{STEC} = \int n_e\,ds$ is the slant total electron content along the line of sight, in units of electrons/m$^2$. One TEC Unit (TECU) equals $10^{16}$\,el/m$^2$.

For a signal from a satellite at zenith angle $z'$ at the ionospheric pierce point (where the line of sight intersects a thin-shell model of the ionosphere at altitude $h_{\text{ion}} \approx 350$\,km), the STEC is related to the vertical TEC (VTEC) by:
\begin{equation}
\text{STEC} = \frac{\text{VTEC}}{\cos z'}.
\label{eq:stec_vtec}
\end{equation}
Combining Eqs.~\eqref{eq:iono_delay} and~\eqref{eq:stec_vtec}, the L1 group delay (at $f_{L1} = 1575.42$\,MHz) is:
\begin{equation}
\Delta I_{L1} = \frac{40.3}{(1575.42 \times 10^6)^2}\,\frac{\text{VTEC}}{\cos z'} \approx \frac{0.162\,\text{VTEC}}{\cos z'}~\text{m},
\label{eq:iono_L1}
\end{equation}
where VTEC is in TECU and $\Delta I_{L1}$ is in meters. For a satellite at zenith ($z' = 0$) and VTEC $= 15$\,TECU, the delay is $\sim\!2.4$\,m; at $z' = 70\degree$ (satellite at $20\degree$ elevation), it increases to $\sim\!7.1$\,m.

\subsubsection{Mapping Ionospheric Delay to Height Bias}
\label{sec:iono_height}

The ionospheric group delay adds a positive range error to every satellite pseudorange: the receiver measures a distance to each satellite that is too long. In the least-squares position solution, this common-mode positive range bias maps \emph{preferentially} into the height (vertical) coordinate, for three related reasons~\cite{coster2008,ovstedal2002}:

\begin{enumerate}[nosep]
\item \textbf{Asymmetric satellite geometry.} GPS satellites are observed only from above (all at positive elevation angles). A symmetric constellation would place some satellites below the horizon, allowing the range bias to cancel. The one-sided geometry means that a positive range bias on all satellites pushes the computed receiver position upward---away from the center of the satellite distribution.

\item \textbf{Elevation-dependent delay.} Low-elevation satellites have longer paths through the ionosphere ($1/\cos z'$ factor in Eq.~\ref{eq:iono_L1}), so they suffer larger delays. The position estimator interprets this as a height-dependent range pattern: the low satellites appear farther away than they should, and the least-squares fit accommodates this by moving the receiver upward.

\item \textbf{Vertical dilution of precision (VDOP).} The GPS constellation provides poorer geometric strength in the vertical than in the horizontal, because no satellites are observed below the horizon. The VDOP is typically 1.5--2.5 (compared to HDOP $\approx$ 0.8--1.5), so range errors are amplified more strongly in the vertical coordinate than in the horizontal.
\end{enumerate}

The combined effect is that the ionospheric height bias is always positive (the receiver appears too high) and is approximately:
\begin{equation}
\delta h_{\text{iono}} \approx \alpha \times 0.162\,\overline{\text{VTEC}}~\text{m},
\label{eq:iono_height_bias}
\end{equation}
where $\alpha \approx 2$--3 is a geometry-dependent amplification factor that accounts for VDOP and the average obliquity of the observed satellites, and $\overline{\text{VTEC}}$ is the effective time-averaged VTEC during the observation period. For $\overline{\text{VTEC}} = 12$\,TECU and $\alpha = 2.5$, this gives $\delta h_{\text{iono}} \approx 4.9$\,m, consistent with the $5.26 \pm 0.93$\,m inferred from the SD data (Eq.~\ref{eq:iono_bias}).

\subsubsection{The Klobuchar Broadcast Model}

The GPS navigation message includes eight ionospheric correction parameters that define the Klobuchar model~\cite{klobuchar1987}. This model approximates the VTEC as a cosine function of local time with a nighttime floor:
\begin{equation}
\Delta I_{\text{Klobuchar}} = F \times
\begin{cases}
5 \times 10^{-9}~\text{s} \times c, & \text{night}, \\[4pt]
\left[5 \times 10^{-9} + A\,\cos\!\left(\dfrac{2\pi(t - 50{,}400)}{P}\right)\right]\!\times c, & \text{day},
\end{cases}
\label{eq:klobuchar}
\end{equation}
where $F$ is the obliquity factor ($\approx 1/\cos z'$), $A$ is the amplitude (from the broadcast parameters, typically 5--50\,ns), $P$ is the period (from the broadcast parameters, typically 72{,}000\,s), $t$ is local time in seconds, and $c$ is the speed of light. The ionosphere is modeled as a thin shell at 350\,km altitude.

The Klobuchar model removes approximately 50\% of the ionospheric delay on an RMS basis~\cite{klobuchar1987}. This is adequate for navigation but leaves a residual bias of several meters that is unacceptable for precise positioning. Critically, if the receiver firmware does not apply even this partial correction---as is believed to be the case for the SD Motorola Oncore receivers---the \emph{full} ionospheric delay biases the position solution.

\subsubsection{Typical TEC Values at the TA Site}

The Telescope Array is located at mid-latitude ($\sim\!39\degree$\,N), a region of moderate ionospheric activity~\cite{mannucci1998}. The VTEC at the TA site follows the expected diurnal and solar-cycle patterns:

\begin{itemize}[nosep]
\item \textbf{Diurnal variation:} VTEC ranges from 5--10\,TECU at night (local midnight through pre-dawn) to 15--30\,TECU during the day (local noon $\pm$ 3\,hours). The peak occurs 1--2 hours after local noon due to the time required for solar heating to maximize the ionization rate.

\item \textbf{Solar cycle variation:} VTEC scales roughly with the solar extreme ultraviolet (EUV) flux, which varies by a factor of 2--3 between solar minimum and solar maximum. The 2011 survey was conducted during solar cycle 24 ascending phase (moderate activity); the SD GPS positions were established over a range of solar conditions spanning 2007--2019.

\item \textbf{Seasonal variation:} At mid-latitudes, VTEC tends to be higher near the equinoxes than the solstices, with amplitude variations of $\sim\!30$\%.

\item \textbf{Geomagnetic storms:} During major storms, VTEC can increase by a factor of 2--3 for hours to days. These events are infrequent ($\sim\!4$--5 major storms per solar cycle) and would produce anomalously large GPS height biases. The $0.93$\,m scatter in the SD $\Delta h$ values may partly reflect storm-time variability.
\end{itemize}

For an ``average'' observation at the TA site, a VTEC of 10--15\,TECU is representative, producing an uncorrected L1 range delay of 1.6--2.4\,m at zenith and a height bias of $\sim\!4$--7\,m after geometric amplification. This range brackets the observed $5.26 \pm 0.93$\,m.


%======================================================================
